\documentclass[a4paper,12pt]{article}
\usepackage[T2A]{fontenc}
\usepackage[utf8]{inputenc}
\usepackage[russian]{babel}
\usepackage{amsmath,amssymb,mathtools}
\usepackage{helvet}
\renewcommand{\familydefault}{\sfdefault}
\usepackage{icomma}
\usepackage[margin=21mm,top=19mm,bottom=19mm]{geometry}
\usepackage{microtype}
\usepackage{tikz}
\usetikzlibrary{arrows.meta,positioning,shapes.geometric,calc}
\usepackage{tabularx,booktabs,array}
\usepackage{enumitem}
\usepackage{hyperref}
\usepackage{parskip}
\usepackage{fancyhdr}
\usepackage{setspace}
\usepackage{xcolor}
\definecolor{accent}{HTML}{007C80}
\definecolor{darkaccent}{HTML}{00575C}
\definecolor{muted}{HTML}{4B555A}
\hypersetup{colorlinks=true,linkcolor=darkaccent,urlcolor=darkaccent,pdfauthor={Лёвин Артём Александрович},pdftitle={Чек-лист. Производительность и совместная работа. 20.11.2026}}
\setstretch{1.08}
\setlength{\parindent}{0pt}
\setlength{\parskip}{4pt}
\setlist[itemize]{leftmargin=6mm,itemsep=3pt,topsep=3pt}
\setlist[enumerate]{leftmargin=7mm,itemsep=9pt,topsep=6pt}
\pagestyle{fancy}
\fancyhf{}
\fancyhead[L]{\footnotesize\textcolor{muted}{Математика / профильный ЕГЭ}}
\fancyhead[R]{\footnotesize\textcolor{muted}{Анна Баннова}}
\fancyfoot[C]{\footnotesize\textcolor{muted}{\thepage}}
\renewcommand{\headrulewidth}{0.25pt}
\renewcommand{\footrulewidth}{0pt}
\setlength{\headheight}{15pt}
\newcommand{\topic}[1]{\section*{\textcolor{darkaccent}{#1}}}
\newcommand{\subtopic}[1]{\subsection*{\textcolor{darkaccent}{#1}}}
\newcommand{\linekey}[1]{\textcolor{darkaccent}{\textbf{#1}}}
\newcommand{\task}[2]{\item \textbf{#1.} #2}
\begin{document}
\begin{titlepage}
\thispagestyle{empty}
\begin{center}
\vspace*{21mm}
{\large\textcolor{accent}{М А Т Е М А Т И К А}}\\[23mm]
{\Huge\bfseries Чек-лист}\\[9mm]
{\LARGE\bfseries Производительность\\[2mm] и совместная работа}\\[11mm]
{\large Текстовые задачи: от таблицы к уравнению}\\[24mm]
{\large Подготовка к ЕГЭ по профильной математике}\\[25mm]
{\Large Анна Баннова}\\[8mm]
{\large 20 ноября 2026 г.}
\vfill
{\normalsize Лёвин Артём Александрович}\\[2mm]
{\normalsize эксклюзивно для Анны Банновой}
\end{center}
\vspace{12mm}
\begin{center}
\begin{tikzpicture}[x=1cm,y=1cm]
\draw[accent!60,line width=0.75pt] (-5,0) .. controls (-2.7,0.8) and (-1.5,-0.55) .. (0,0) .. controls (1.7,0.62) and (3.1,-0.70) .. (5,0.12);
\end{tikzpicture}
\end{center}
\end{titlepage}

\topic{1. Как узнавать задачу на производительность}
\linekey{Сигнал в условии:} «литров в минуту», «деталей в час», «страниц за минуту», «изделий в день». Такие величины описывают \textbf{скорость выполнения работы}, или производительность.

\begin{center}
\begin{tabularx}{\linewidth}{@{}>{\bfseries}p{0.28\linewidth}X X@{}}
\toprule
Величина & Что означает & Единицы измерения \\
\midrule
Результат $A$ & Сколько всего сделано & л, деталей, страниц, кг \\
Время $t$ & Сколько длилась работа & мин, ч, дни \\
Производительность $p$ & Сколько делают за единицу времени & л/мин, деталей/ч и т.~д. \\
\bottomrule
\end{tabularx}
\end{center}

\begin{equation*}
\boxed{A=p\,t}\qquad \boxed{p=\frac{A}{t}}\qquad \boxed{t=\frac{A}{p}}.
\end{equation*}

\linekey{Универсальная таблица} начинается с трёх колонок:
\begin{center}
\begin{tabularx}{\linewidth}{@{}lXXX@{}}
\toprule
Объект & Производительность $p$ & Время $t$ & Результат $A$ \\
\midrule
Первый & \dots & \dots & \dots \\
Второй & \dots & \dots & \dots \\
Вместе & \dots & \dots & \dots \\
\bottomrule
\end{tabularx}
\end{center}

\textbf{Как заполнить таблицу:}
\begin{enumerate}[itemsep=2pt]
\item Прочитайте, \textbf{что именно требуется найти}. Если нужны литры в минуту, неизвестной обычно удобно сделать производительность. Если нужны минуты, неизвестной удобно сделать время.
\item Запишите данные про каждую трубу, насос, станок или рабочего. Не смешивайте единицы измерения.
\item Получите недостающие клетки с помощью формул $A=pt$, $p=A/t$, $t=A/p$.
\item Переведите сравнение или совместную работу в уравнение.
\end{enumerate}

\topic{2. Какое уравнение составить}
\textbf{Сравнение по времени.} Если первый объект работает на $d$ минут \emph{дольше} второго, то
\begin{equation*}
\boxed{t_{\text{первого}}-t_{\text{второго}}=d}.
\end{equation*}
\textbf{Совместная работа.} Если два объекта выполняют одну работу одновременно и их производительности постоянны, то
\begin{equation*}
\boxed{p_{\text{вместе}}=p_1+p_2},\qquad
\boxed{t_{\text{вместе}}=\frac{A_{\text{общий}}}{p_1+p_2}}.
\end{equation*}
\textcolor{muted}{Оговорка: складывать производительности можно тогда, когда оба объекта действительно работают одновременно, а выполненные объёмы суммируются.}

\clearpage
\topic{3. Разобранный пример: две трубы}
\textbf{Условие.} Первая труба пропускает на $12$ литров воды в минуту меньше, чем вторая. Каждая труба по отдельности может заполнить бак объёмом $160$ литров. Первой трубе для этого требуется на $12$ минут больше, чем второй. Найдите производительность первой трубы.

Пусть $x$~л/мин --- производительность первой трубы. Тогда производительность второй $x+12$~л/мин, причём $x>0$.
\begin{center}
\begin{tabularx}{\linewidth}{@{}lXXX@{}}
\toprule
Труба & Производительность & Время & Результат \\
\midrule
Первая & $x$ & $\dfrac{160}{x}$ & $160$ л \\
Вторая & $x+12$ & $\dfrac{160}{x+12}$ & $160$ л \\
\bottomrule
\end{tabularx}
\end{center}

Первая труба медленнее, поэтому работает дольше:
\[
\frac{160}{x}-\frac{160}{x+12}=12.
\]
Чтобы облегчить вычисления, \textbf{делим все слагаемые уравнения на $4$}:
\[
\frac{40}{x}-\frac{40}{x+12}=3.
\]
Приводим дроби к общему знаменателю, сохраняя скобки у вычитаемого:
\begin{align*}
\frac{40(x+12)-40x}{x(x+12)}&=3,\\
\frac{40x+480-40x}{x^2+12x}&=3,\\
\frac{480}{x^2+12x}&=\frac{3}{1}.
\end{align*}
\textbf{Перемножаем крест-накрест} (знаменатель положителен при $x>0$):
\[
3(x^2+12x)=480\quad\Longrightarrow\quad x^2+12x-160=0.
\]
\[
D=12^2-4\cdot1\cdot(-160)=144+640=784=28^2.
\]
\[
x_1=\frac{-12+28}{2}=8,\qquad x_2=\frac{-12-28}{2}=-20.
\]
Отрицательный корень не подходит: производительность должна быть положительной.

\linekey{Проверка:} первая труба перекачивает $8$~л/мин, вторая $20$~л/мин. Время: $160/8=20$~мин и $160/20=8$~мин. Разница действительно $12$~мин.

\linekey{Ответ: $8$ л/мин.}

\clearpage
\topic{4. Разобранный пример: работа двух насосов вместе}
\textbf{Условие.} Первый садовый насос перекачивает $11$ литров воды за $6$ минут, второй --- такой же объём за $7$ минут. Оба насоса включили одновременно. За сколько минут они совместно перекачают $143$ литра воды?

Пусть $x$ --- искомое время совместной работы, $x>0$.
\begin{center}
\begin{tabularx}{\linewidth}{@{}lXXX@{}}
\toprule
Режим & Производительность & Время & Результат \\
\midrule
Насос 1 & $\dfrac{11}{6}$ л/мин & $6$ мин & $11$ л \\
Насос 2 & $\dfrac{11}{7}$ л/мин & $7$ мин & $11$ л \\
Оба насоса & $\dfrac{143}{x}$ л/мин & $x$ мин & $143$ л \\
\bottomrule
\end{tabularx}
\end{center}

Суммарная производительность равна сумме отдельных:
\begin{align*}
\frac{143}{x}&=\frac{11}{6}+\frac{11}{7}\\
&=\frac{11\cdot7+11\cdot6}{6\cdot7}\\
&=\frac{11(7+6)}{6\cdot7}=\frac{11\cdot13}{42}=\frac{143}{42}.
\end{align*}
Теперь $143/x=143/42$, поэтому $x=42$.

\linekey{Проверка:} вместе насосы перекачивают $143/42$ литра за минуту. За $42$ минуты получается ровно $143$ литра.

\linekey{Ответ: $42$ минуты.}

\subtopic{Три приёма, которые экономят вычисления}
\begin{enumerate}[itemsep=3pt]
\item \textbf{«Улыбочка» для дробей:} $\displaystyle\frac{a}{b}\pm\frac{c}{d}=\frac{ad\pm bc}{bd}$ при $b\ne0$, $d\ne0$. Вычитаемую группу в длинном числителе лучше заключать в скобки.
\item \textbf{Общий множитель:} $11\cdot7+11\cdot6=11(7+6)$. Часто это сохраняет небольшие числа до конца решения.
\item \textbf{Пропорция:} из $\displaystyle\frac{a}{b}=\frac{c}{d}$ при $b,d\ne0$ следует $ad=bc$. Переход выполняют только после проверки допустимости знаменателей.
\end{enumerate}

\textbf{Дополнительная тренировка преобразования.} Если
\[
\frac{100}{x}-\frac{100}{x+5}=1,\quad x>0,
\]
то $\frac{500}{x(x+5)}=1$, а значит $x^2+5x-500=0$. Здесь $D=2025=45^2$, допустимый корень $x=(-5+45)/2=20$.

\clearpage
\topic{5. Квадратные корни без калькулятора}
При решении текстовых задач часто появляется квадратное уравнение $ax^2+bx+c=0$. Его дискриминант
\[
D=b^2-4ac,\qquad x=\frac{-b\pm\sqrt D}{2a}\quad(a\ne0,\ D\ge0).
\]

\subtopic{Как быстро найти корень из полного квадрата}
\begin{enumerate}[itemsep=5pt]
\item \textbf{Найдите десяток.} Например, $400<784<900$, поэтому $20<\sqrt{784}<30$.
\item \textbf{Посмотрите на последнюю цифру.} Число $784$ оканчивается на $4$, поэтому целый корень может оканчиваться на $2$ или $8$. Кандидаты в данном диапазоне: $22$ и $28$.
\item \textbf{Проверьте кандидата возведением в квадрат.} $28^2=(30-2)^2=900-120+4=784$. Значит, $\sqrt{784}=28$.
\end{enumerate}

\begin{center}
\begin{tabular}{@{}lcl@{}}
\toprule
Число & Кандидаты & Проверенный корень \\
\midrule
$5184$ & $72$ или $78$ & $72^2=5184$ \\
$7921$ & $81$ или $89$ & $89^2=7921$ \\
$3249$ & $53$ или $57$ & $57^2=3249$ \\
$1764$ & $42$ или $48$ & $42^2=1764$ \\
$2025$ & $45$ & $45^2=2025$ \\
\bottomrule
\end{tabular}
\end{center}

\linekey{Полезные ориентиры:} $10^2=100$, $20^2=400$, $30^2=900$, $40^2=1600$, $50^2=2500$, $60^2=3600$, $70^2=4900$, $80^2=6400$, $90^2=8100$, $100^2=10000$.

\textbf{Важная граница метода:} последняя цифра даёт только кандидатов. Близость числа к соседнему круглому квадрату помогает выбрать наиболее вероятный вариант, но правильность \textbf{обязательно подтверждается возведением в квадрат}. Если число не является полным квадратом, целый корень искать бессмысленно.

\subtopic{Самопроверка перед ответом}
\begin{itemize}
\item Все времена измерены в одинаковых единицах, а объёмы --- в одинаковых единицах?
\item Если объект медленнее, действительно ли он тратит больше времени на \emph{одинаковую} работу?
\item В строке совместной работы сложены производительности, а не произвольные числа из условия?
\item Перед делением и перемножением крест-накрест проверены знаменатели?
\item При раскрытии скобок каждое слагаемое умножено на внешний множитель?
\item Проверено, подходит ли найденный корень по смыслу и подстановке?
\end{itemize}

\clearpage
\thispagestyle{fancy}
\begin{center}
{\Large\bfseries\textcolor{darkaccent}{Алгоритм: как решать задачу на производительность}}
\end{center}
\vspace{6mm}
\begin{center}
\tikzset{
  algstep/.style={draw=accent!75,rounded corners=2pt,align=center,text width=11.6cm,minimum height=1.22cm,inner sep=6pt,fill=white},
  algdecision/.style={diamond,draw=accent!75,aspect=3.0,align=center,text width=5.2cm,inner sep=3pt,fill=white},
  algarrow/.style={-{Latex[length=3mm]},line width=0.8pt,color=darkaccent}
}
\begin{tikzpicture}[node distance=11mm and 28mm, every node/.style={font=\small}]
\node[algstep] (read) {Прочитайте условие. Установите, что нужно найти:\\результат, время или производительность.};
\node[algstep,below=of read] (table) {Составьте таблицу «производительность --- время --- результат».\\Обозначьте искомое за $x$ и заполните строки.};
\node[algdecision,below=12mm of table] (kind) {Как связаны объекты?};
\node[algstep,text width=5.1cm,below left=15mm and 4mm of kind] (compare) {Одинаковая работа,\\разное время\\$t_{\text{дольше}}-t_{\text{быстрее}}=d$};
\node[algstep,text width=5.1cm,below right=15mm and 4mm of kind] (joint) {Работают одновременно\\$p_{\text{вместе}}=p_1+p_2$};
\node[algstep,below=63mm of kind] (algebra) {Получите уравнение. Сокращайте общий множитель,\\приводите дроби к знаменателю и решайте.};
\node[algstep,below=of algebra] (check) {Проверьте ОДЗ, единицы, знак корня и подстановку в условие.\\Запишите ответ с единицами измерения.};
\draw[algarrow] (read) -- (table);
\draw[algarrow] (table) -- (kind);
\draw[algarrow] (kind.south west) -| node[pos=0.25,above left,font=\scriptsize]{Сравнение} (compare.north);
\draw[algarrow] (kind.south east) -| node[pos=0.25,above right,font=\scriptsize]{Сумма} (joint.north);
\draw[algarrow] (compare.south) -- ++(0,-6mm) -| ($(algebra.north)+(-2cm,0)$);
\draw[algarrow] (joint.south) -- ++(0,-6mm) -| ($(algebra.north)+(2cm,0)$);
\draw[algarrow] (algebra) -- (check);
\end{tikzpicture}
\end{center}

\clearpage
\topic{6. Тренировочный блок --- 10 задач}
\textcolor{muted}{Идите по порядку: от сравнения времени к совместной работе. Для каждой текстовой задачи составьте таблицу; затем запишите уравнение, решите его и проверьте ответ по смыслу. Не забывайте указывать единицы измерения.}

\subtopic{Уровень 1. Сравнение времени и производительности}
\begin{enumerate}[label=\textbf{\arabic*.},leftmargin=9mm,itemsep=11pt]
\task{Две водопроводные трубы}{В мастерской установлен резервуар объёмом $120$ литров. Первая труба пропускает на $4$ литра в минуту меньше второй. Если открывать трубы по отдельности и каждый раз заполнять резервуар с нуля, первой трубе потребуется на $5$ минут больше. Сколько литров воды в минуту пропускает \textbf{первая} труба?}
\task{Два принтера}{Типографии нужно напечатать одинаковые партии по $180$ рекламных листовок. Первый принтер медленнее второго на $6$ листовок в минуту. На одну такую партию первый принтер тратит на $5$ минут больше второго. Сколько листовок за минуту печатает \textbf{второй} принтер?}
\task{Упаковочные линии}{На складе проверяют две упаковочные линии. Каждая отдельно должна упаковать партию из $360$ одинаковых коробок. Быстрая линия упаковывает на $10$ коробок в минуту больше медленной и заканчивает работу на $3$ минуты раньше. За сколько минут выполнит партию \textbf{медленная} линия?}
\task{Шланги для полива}{На садовом участке есть два шланга, которыми можно по отдельности наполнить ёмкость для полива объёмом $240$ литров. Второй шланг подаёт на $8$ литров в минуту больше первого, поэтому наполняет ёмкость на $8$ минут быстрее. Найдите производительность \textbf{второго} шланга в литрах в минуту.}
\task{Печать открыток}{Две машины по очереди изготовляют одинаковые партии по $336$ открыток. Более быстрая машина делает на $4$ открытки в минуту больше медленной. Медленной машине на изготовление одной партии требуется на $7$ минут больше. Сколько открыток в минуту изготавливает \textbf{медленная} машина?}
\end{enumerate}

\clearpage
\subtopic{Уровень 2. Суммарная производительность}
\begin{enumerate}[label=\textbf{\arabic*.},start=6,leftmargin=9mm,itemsep=13pt]
\task{Откачка воды из подвала}{Первый насос за $5$ минут откачивает $50$ литров воды, а второй за $7$ минут откачивает $42$ литра. Скорость работы каждого насоса постоянна. Насосы включили одновременно, чтобы удалить из подвала $320$ литров воды. Через сколько минут после включения они вместе закончат работу?}
\task{Два фасовочных станка}{На производстве первый станок фасует $24$ килограмма крупы за $3$ минуты, а второй --- $36$ килограммов за $6$ минут. Весь объём крупы можно разделить между станками, оба работают непрерывно с постоянной производительностью. За сколько минут совместной работы они расфасуют $280$ килограммов крупы?}
\task{Сортировка посылок}{На логистической площадке первая сортировочная линия обрабатывает $28$ посылок за $4$ минуты, а вторая --- $27$ посылок за $3$ минуты. Обе линии могут работать одновременно, и каждая сохраняет указанную скорость. Сколько минут понадобится, чтобы совместно обработать $352$ посылки?}
\task{Два насоса в оранжерее}{В оранжерее первый насос перекачивает $44$ литра питательного раствора за $4$ минуты, а второй --- $60$ литров за $5$ минут. Для наполнения большого резервуара их запускают одновременно; суммарно требуется подать $529$ литров. Найдите время наполнения резервуара в минутах.}
\end{enumerate}

\subtopic{Уровень 3. Работа в два этапа}
\begin{enumerate}[label=\textbf{\arabic*.},start=10,leftmargin=9mm,itemsep=9pt]
\task{Заполнение накопительного бака}{Бак объёмом $580$ литров наполняют двумя насосами. Первый насос подаёт $16$ литров воды в минуту и начинает работу один. Через $10$ минут после его включения запускают второй насос, который подаёт $12$ литров в минуту. Далее оба работают одновременно до полного наполнения бака. Сколько \textbf{минут прошло с момента включения первого насоса} до заполнения бака?}
\end{enumerate}

\clearpage
\topic{7. Ответы и контрольная проверка}
\small
\renewcommand{\arraystretch}{1.45}
\begin{tabularx}{\linewidth}{@{}>{\bfseries}p{0.07\linewidth}p{0.23\linewidth}X@{}}
\toprule
№ & Ответ & Краткая проверка \\
\midrule
1 & $8$ л/мин & $120/8-120/12=15-10=5$ мин \\
2 & $18$ листовок/мин & $180/12-180/18=15-10=5$ мин \\
3 & $12$ мин & $360/30-360/40=12-9=3$ мин \\
4 & $20$ л/мин & $240/12-240/20=20-12=8$ мин \\
5 & $12$ открыток/мин & $336/12-336/16=28-21=7$ мин \\
6 & $20$ мин & $50/5+42/7=10+6=16$ л/мин; $320/16=20$ \\
7 & $20$ мин & $24/3+36/6=8+6=14$ кг/мин; $280/14=20$ \\
8 & $22$ мин & $28/4+27/3=7+9=16$ посылок/мин; $352/16=22$ \\
9 & $23$ мин & $44/4+60/5=11+12=23$ л/мин; $529/23=23$ \\
10 & $25$ мин & $16\cdot10+(16+12)\cdot15=160+420=580$ л \\
\bottomrule
\end{tabularx}
\normalsize

\vspace{11mm}
\subtopic{Итоговая памятка}
\begin{tabularx}{\linewidth}{@{}p{0.41\linewidth}X@{}}
\toprule
Сигнал в условии & Полезное действие \\
\midrule
«На $d$ минут дольше» & Из большего времени вычесть меньшее и приравнять разность к $d$ \\
«Вместе», «одновременно» & Сложить производительности: $p=p_1+p_2$ \\
«Сколько литров за минуту?» & Неизвестная величина --- производительность \\
«Сколько минут потребуется?» & Неизвестная величина --- время \\
В выражении много больших чисел & Поискать общий множитель или общий делитель всех частей уравнения \\
Разность алгебраических дробей & Привести к общему знаменателю; при раскрытии скобок сохранить знак минус \\
Получился корень из дискриминанта & Ограничить диапазон, найти возможную последнюю цифру, проверить квадратом \\
\bottomrule
\end{tabularx}

\vfill
\begin{center}\textcolor{muted}{\footnotesize Лёвин Артём Александрович --- эксклюзивно для Анны Банновой}\end{center}
\end{document}
